%! AP Statistics — Unit 2, Lesson 2.3 — Warm-Up
\documentclass[10pt]{article}
\usepackage{apstats-article}
\usepackage{apstats-boxes}

%=============================================================================
\begin{document}

\pageheader{Unit 2, Lesson 2.3}{Warm-Up}
\namedateperiod

\vspace{0.3in}

% ── SECTION 1: Spiral Review ──────────────────────────────────────────────────
{\large\bfseries\color{navy} Part 1 \quad Conditional Relative Frequencies (Lesson 2.2 Review)}

\vspace{0.12in}

\noindent A survey of 200 high school students recorded two variables: whether they play
a sport (Yes / No) and whether they have a part-time job (Yes / No). Results are below.

\vspace{0.1in}

\renewcommand{\arraystretch}{1.9}
\begin{tabularx}{\linewidth}{@{} l C{2.8cm} C{2.8cm} C{2.4cm} @{}}
\rowcolor{sky}
                    & \textbf{Job: Yes} & \textbf{Job: No} & \textbf{Row Total} \\
\hline
\textbf{Sport: Yes} & 30 & 90 & \blank{1.8cm} \\
\textbf{Sport: No}  & 40 & 40 & \blank{1.8cm} \\
\hline
\textbf{Col.\ Total}& \blank{1.8cm} & \blank{1.8cm} & \blank{1.8cm} \\
\end{tabularx}

\vspace{0.18in}

\begin{enumerate}

    \item Fill in all missing row totals, column totals, and the grand total.

    \vspace{0.06in}

    \item Compute the conditional relative frequency of having a job for student-athletes.\\[8pt]
    $P(\text{Job: Yes} \mid \text{Sport: Yes}) = \dfrac{\blank{1.2cm}}{\blank{1.2cm}} = \blank{2.2cm}$

    \vspace{0.18in}

    \item Compute the conditional relative frequency of having a job for non-athletes.\\[8pt]
    $P(\text{Job: Yes} \mid \text{Sport: No}) = \dfrac{\blank{1.2cm}}{\blank{1.2cm}} = \blank{2.2cm}$

    \vspace{0.14in}

    \item Based on your answers to \#2 and \#3, does there appear to be an association between
    playing a sport and having a job? Circle one and explain briefly.\\[4pt]
    \textbf{Yes} \quad / \quad \textbf{No}\\[8pt]
    \writeline\\[2pt]\writeline

\end{enumerate}

\vspace{0.28in}

% ── SECTION 2: Looking Ahead ──────────────────────────────────────────────────
{\large\bfseries\color{navy} Part 2 \quad Putting a Number on the Difference}

\vspace{0.12in}

\begin{tcolorbox}[colback=greenbg, colframe=greenacc, boxrule=0.8pt, arc=2mm,
    left=4mm, right=4mm, top=3mm, bottom=3mm]
    \small
    \textbf{Think about it:} In Part 1 you found two conditional proportions and compared
    them informally. Today we will compute \emph{statistics} that quantify exactly
    \emph{how different} the two proportions are.

    \vspace{6pt}
    Two questions to consider before the lesson:
    \begin{itemize}[leftmargin=1.3em, itemsep=6pt]
        \item \textbf{Additive:} How many percentage points apart are the two proportions?
              This is the \textbf{difference in proportions}.
        \item \textbf{Multiplicative:} How many times as likely is one group compared to the other?
              This is the \textbf{relative risk}.
    \end{itemize}
\end{tcolorbox}

\vspace{0.14in}

\noindent Using your values from Part 1, \textbf{subtract} the two conditional proportions.\\[8pt]
$\hat{p}_{\text{athlete}} - \hat{p}_{\text{non-athlete}} = \blank{2cm} - \blank{2cm} = \blank{2.2cm}$

\vspace{0.1in}

\noindent What does this value tell you? Write one sentence.\\[8pt]
\writeline\\[2pt]\writeline

\end{document}
